% Main-text candidate. Environments proposition/theorem/corollary/remark required.
% All identifiers start sc: to avoid collisions with matrix-agent contributions.

\subsection{From one deterministic edge to several stochastic diagnostics}
On the scalar sampled quadratic
\(\ell_B(x)=h_Bx^2/2-\xi_Bx\), SGD is exactly
\begin{equation}
 X_{t+1}=A_{t+1}X_t+\eta\xi_{t+1},\qquad A_B=1-\eta h_B.
 \label{sc:quadratic}
\end{equation}
Two trajectories using the same sampled batches differ by
\(\Delta_{t+1}=A_{t+1}\Delta_t\).
Write \(M=|A_B|\), so moments always use the absolute multiplier, and define
\begin{equation}
 m(p)=\mathbb E M^p,\quad
 \gamma=\mathbb E\log M,\quad
 H=\frac1{m(-1)}.
 \label{sc:diagnostics}
\end{equation}
Here \(H\) is the harmonic gain; its inverse \(m(-1)\) is the negative-moment
coefficient. We only assign a finite inverse-moment estimate this meaning
when the population inverse moment exists. For constant curvature,
\(m(p)^{1/p}=M\), \(\gamma=\log M\), and \(H=M\): the boundaries all
coincide at \(|1-\eta h|=1\). Curvature randomness separates them.

\begin{proposition}[Scalar moment ordering]\label{sc:ordering}
Let \(M>0\), \(\mathbb E|\log M|<\infty\), and suppose the displayed moments
are finite. Then
\begin{equation}
 \log H\ \le\ \gamma\ \le\ \log m(1)\ \le\ \tfrac12\log m(2).
 \label{sc:ladder}
\end{equation}
For iid multipliers and \(\Delta_0\ne0\),
\(t^{-1}\log|\Delta_t/\Delta_0|\to\gamma\) almost surely and
\(\mathbb E|\Delta_t/\Delta_0|^p=m(p)^t\).
Consequently \(H>1\) implies \(\gamma>0\). The converse need not hold.
\end{proposition}

The classical Kesten--Goldie result concerns the \emph{affine} recursion
\eqref{sc:quadratic}: with \(\gamma<0\), nondegenerate additive forcing, and
the usual renewal assumptions, its stationary upper tail has exponent
\(\kappa>0\) solving \(m(\kappa)=1\)
\citep{kesten1973random,goldie1991implicit}.
This is a baseline, not a theorem of post-expansion confinement. In
particular, the nondegenerate Gaussian affine model used in
Appendix~\ref{sc:proofs} has no invariant probability when \(\gamma>0\).
An additional return mechanism is needed.

\subsection{Nonlinear return changes an upper-tail question into a central-mass question}
The simplest exact return model is saturated amplitude:
\begin{equation}
 R_{t+1}=\min\{r_*,M_{t+1}R_t\},\qquad 0<R_t\le r_*.
 \label{sc:saturation}
\end{equation}
It separates local multiplicative motion from return at the outer scale.
It is an explicit model, not an identity for a globally quadratic loss.

\begin{theorem}[An inverse-moment law after local expansion]\label{sc:central-law}
Suppose \(\mathbb E|\log M|<\infty\), \(\gamma>0\),
\(\log M\) is nonarithmetic, and
\(m(-\alpha)=1\) for some \(\alpha>0\). Assume \(m(-q)\) is finite in a
neighborhood of \(\alpha\) and
\(0<-\mathbb E[M^{-\alpha}\log M]<\infty\).
Then \eqref{sc:saturation} has a unique invariant probability on
\((0,r_*]\), and for some \(C\in(0,\infty)\),
\begin{equation}
 \mathbb P_\pi(R\le r)\sim C(r/r_*)^\alpha,\qquad r\downarrow0.
 \label{sc:smallball}
\end{equation}
If \(\log M\) additionally has a bounded continuous density with compact
support, its radial density satisfies
\[
 p_R(r)\sim C\alpha r_*^{-\alpha}r^{\alpha-1}.
\]
When the negative-moment curve is strictly convex and finite through
\(1\) and \(\alpha\),
\begin{equation}
 H<1\Longleftrightarrow\alpha<1,\qquad
 H=1\Longleftrightarrow\alpha=1,\qquad
 H>1\Longleftrightarrow\alpha>1.
 \label{sc:shape-ladder}
\end{equation}
\end{theorem}

Thus positive local log growth can coexist with a central density cusp
\((0<\alpha<1)\); the harmonic crossing separates this cusp from central
depletion \((\alpha>1)\) in this return model. For a reflection-symmetric
signed coordinate, the same power occurs on each side. Central depletion
is a local shape conclusion. Exactly two modes additionally requires
control of the outer density, and arbitrary reinjection can reverse the
conclusion (Appendix~\ref{sc:counterexample}).

\paragraph{Proof idea.}
The change of variable \(W=\log(r_*/R)\) gives
\(W_{t+1}=\max\{0,W_t-\log M_{t+1}\}\).
This is a reflected random walk with negative mean increment.
Its invariant law is the supremum of the unreflected walk. Tilting the
increment law by \(M^{-\alpha}\) reverses its drift; the overshoot renewal
theorem gives an exponential tail in \(W\), hence the power
\eqref{sc:smallball} in \(R\). The root therefore has the opposite sign
from the Kesten upper-tail root.

\subsection{An exact smooth SGD construction beyond both central boundaries}
To prove that return can arise from a gradient update, rather than from
an imposed cap, draw \(h\in\{.95,1.05\}\) equiprobably and
\(e\sim N(0,\sigma^2)\), independently, with \(\sigma>0\). Use
\begin{equation}
 \ell(x;h,e)=\tfrac12(\sqrt h\,x-e)^2
 -\tfrac{.92}{2}\{x^2-\log(1+x^2)\}.
 \label{sc:smooth-loss}
\end{equation}
Every sampled loss is coercive. The population loss, up to a constant, is
\(L(x)=.04x^2+.46\log(1+x^2)\), with unique minimizer zero but negative
curvature on part of the line.

\begin{theorem}[Smooth nonconvex SGD can return despite central expansion]
\label{sc:nonlinear-recurrence}
More generally, replace \(.92\) in \eqref{sc:smooth-loss} by \(\beta\)
and draw \(h=\beta+\mu\pm s\) equiprobably, with
\(0<s<\mu\) and \(\beta>8\mu\). At batch size one, every step
\(0<\eta<2/(\mu+s)\) has a unique attracting invariant probability
with a smooth, strictly positive, even density and all finite polynomial
moments. Its central Jacobian is \(1-\eta h\), whereas its outer slopes
are \(1-\eta(\mu\pm s)\).

For \((\beta,\mu,s,\eta)=(.92,.08,.05,2.015)\), the central absolute
gains are \(.91425,1.11575\), and
\begin{equation}
 \gamma_0=0.009938\ldots>0,\qquad
 H_0=1.004999\ldots>1,\qquad \alpha_0=2.017427\ldots>1,
 \label{sc:concrete-numbers}
\end{equation}
where \(\tfrac12(.91425^{-\alpha_0}+1.11575^{-\alpha_0})=1\).
The outer slopes are \(.93955,.73805\), both strictly contracting.
\end{theorem}

\paragraph{Proof idea.}
The exact update is
\(X^+=A_\infty X-\eta\beta X/(1+X^2)+\eta\sqrt h\,e\), with
\(A_\infty=1-\eta(h-\beta)\). The middle term is bounded and the outer
slopes contract. A polynomial Foster drift and the positive Gaussian
transition density give recurrence and uniqueness. The local derivative
at the origin is expansive because the sampled curvature is much larger
there. This theorem proves recurrence after the \emph{central} Lyapunov
and harmonic boundaries; it does not identify the tangent exponent along
the invariant trajectory or infer its modes. These are separate empirical
measurements below.
